\section{The observable determines the topology boundary}
\label{sec:main-observables}

Replacing finite resistance sets by their interval hulls is attractive,
since continuous coefficients can be eliminated through convex images,
whereas a product of finite sets has exponentially many choices. Whether
the replacement preserves an extremum depends on the graph \emph{and}
the objective. In this section a universal statement
quantifies over every fixed balanced nomination and all independent
compact positive resistance sets with attained endpoints, under quadratic
laws and without operating filters.

\begin{theorem}[Universal resistance-hull and region classifications]
\label{thm:main-hulls}
For finite connected simple graphs, the universal interval-hull classes
for maxima and minima are exactly those in \cref{tab:main-hulls}.
The same classes characterize separate monotonicity in every resistance
coordinate, with the other coordinates fixed. The attainable flow region
is convex for every fixed nomination and positive resistance box exactly
on cacti. Normalized potential regions and joint flow--potential regions
are universally convex exactly on trees.
\end{theorem}

\begin{table}[tb]
\centering\small
\begin{tabularx}{\linewidth}{@{}Yl@{}}
\toprule
Observable whose extrema must be preserved & Universal graph class\\
\midrule
Every balanced weighted potential $c^\top\pi$ & Trees\\
Every prescribed difference $\pi_s-\pi_t$ & Cacti\\
Every linear flow objective $d^\top x$ & Cacti\\
Every individual arc flow $x_a$ & Series-parallel\\
\bottomrule
\end{tabularx}
\caption{Exact equality of resistance-set and interval-hull extrema
under fixed nominations, independent compact positive quadratic
resistance sets, and no operating filters. Equality of extrema does not
assert an exact scalar comparison algorithm.}
\label{tab:main-hulls}
\end{table}

These claims are \cref{thm:a-weight-hierarchy,thm:a-weight-convex-regions},
with the individual-arc equivalence proved in
\cref{thm:a-bound-characterization}. Counterexamples use rational data,
one varying resistance, and strictly positive resistances on all other
edges. The full-region
classifications extend qualitatively to every fixed common power
$\sgn(x)|x|^q$ with $q>0$, $q\ne1$, without a uniform bit bound or
algorithm (\cref{thm:a-weight-power-regions}). With linear laws, by
contrast, every graph has the hull property for potential objectives
(\cref{sec:a-weighted}).

\subsection{Finite comparison and deterministic envelopes}

A finite secant argument explains the positive individual-arc result
without differentiating a law at zero. The difference of two passive
states for the same nomination and two law profiles is an electrical
network with positive secant resistances and explicit edge disturbances.
For a target edge $a=(u,v)$, let $j$ be the unit electrical current from
$u$ to $v$ in that secant network. The target flow difference is a
linear combination of the disturbances, with coefficients $j_e$ for
$e\ne a$ and $-(1-j_a)$ on $a$, divided by the positive target secant
resistance. By classical confluence, the signs of these currents do not
depend on the positive resistances when the graph is series-parallel.

One may therefore select, pointwise in flow, a maximum or minimum
allowed law on each edge according to the sign of its current $j_e$. The target
state of the resulting deterministic envelope attains the uncertainty
extremum, and compactness makes the selected law value realizable by an
original parameter at that state. The construction is uniform in the
nomination, although the realizing profile can depend on the state
(\cref{thm:a-bound-envelope-structural}, for independent compact
continuous law families). For quadratic interval or finite resistance
uncertainty the envelope is directional quadratic, and its primitive
has a linear-size second-order-cone formulation. Rational convex
optimization and recovery then give polynomial additive arc optimization
for fixed nominations on every series-parallel graph, without a rank
bound (\cref{thm:a-bound-envelope-algorithm}).

Deterministic feasibility for rated monotone characteristics has related
series/parallel propagation methods: Zemanian's report uses endpoint
propagation and operating-point recovery
\citep[Conditions 2.1, Lemmas 4.1 and 5.1, Procedure 8.1]{zemanian1999-ratings}.
Those deterministic procedures do not provide the present
uncertainty-envelope and accuracy-bit output guarantees.

\subsection{What fails when the objective changes}

A triangle already separates weighted potentials from pairwise
potentials. With edges oriented $0\to1\to2\to0$, resistances
$(1,1,\theta)$, nominations $(2,-3,1)$, and
$c=(-5,12,-7)$, the circulation $q$ satisfies
$6q+3-\theta q^2=0$ on $-1/2<q<0$. Its weighted potential is
\[
 W=-\frac{105}{4}-12(q+1/4)^2.
\]
At $\theta=24$, $q=-1/4$ and $W=-105/4$. At
$\theta=16/3$ and $144$, the circulation is respectively $-3/8$ and
$-1/8$ and both values are smaller by $3/16$. Thus a two-point set
and its interval hull give different weighted-potential optima on one
cycle, although this graph is in the universal class for every
individual potential difference. Weighted cancellation changes the
extremal geometry.

A rank-two theta block supplies the corresponding linear-flow
counterexample; a rank-three $K_4$ supplies the individual-arc
counterexample. Subdividing suitable uncertain paths encodes Subset
Sum at the same small block rank, giving the discrete boundaries in
\cref{tab:main-boundaries}, with explicit rational gaps and positive
integer resistances where stated. Multiplying all resistances scales a
pressure gap but leaves flows unchanged; flow gaps are amplified by
scaling nominations instead.

\begin{table}[tb]
\centering\small
\begin{tabularx}{\linewidth}{@{}P{.22\linewidth}P{.26\linewidth}Y@{}}
\toprule
Objective and graph & Scenario assumptions & Computational conclusion\\
\midrule
Potential difference; one rank-two block & Fixed small nominations;
independent two-point resistances & NP-hard, with polynomial-bit gap
(\cref{thm:a-bound-pressure}).\\[3pt]
Individual arc; series-parallel, fixed block rank & Balanced nomination
box; independent finite resistances & Exact algebraic extrema and robust
capacity decision; includes every rank-at-most-two graph
(\cref{cor:a-bound-finite-positive}).\\[3pt]
Individual arc; rank-three block & Fixed small nominations;
independent two-point resistances & Existential strict or weak threshold
is NP-complete; complementary universal test is coNP-complete
(\cref{thm:a-bound-arc-hardness}).\\[3pt]
Weighted potential; one cycle & Fixed nominations and three objective
support vertices; independent finite resistances & NP-complete threshold
attainment, coNP-complete robust upper bounds
(\cref{thm:a-weight-cycle-hard}).\\[3pt]
Linear flow; rank-two block & Fixed nominations; independent finite
resistances & NP-complete threshold attainment; includes unit-cost total
absolute flow on a DAG (\cref{thm:a-weight-flow-hard}).\\[3pt]
Individual arc; series-parallel with growing block rank & Balanced
nomination box; fixed resistances & NP-hard optimization and coNP-hard
robust upper bounds; no general membership claim
(\cref{thm:a-bound-nomination}).\\
\bottomrule
\end{tabularx}
\caption{Discrete and nomination boundaries for quadratic laws. All
rows use unfiltered scenarios. Hardness gaps obstruct polynomial
accuracy-bit algorithms unless P equals NP; they are not uniformly
strong-hardness claims.}
\label{tab:main-boundaries}
\end{table}

Two further distinctions matter. First, the envelope theorem and hull
equality do not settle exact scalar arithmetic on unrestricted
series-parallel graphs. A parallel probe on a cactus branch transfers its radical sum
to an arc comparison, giving Square-Root-Sum hardness even with fixed
unit nominations and fixed integral resistances; a restricted probe
family has the converse reduction (\cref{thm:a-bound-arc-srs}). Second,
hull preservation for an extremum does not preserve exact realization
of a target flow. For a prescribed rational $\bar x$, conservation can
be checked and the equations $A^\top\pi=(\beta_e\bar x_e|\bar x_e|)_e$
are linear in $\pi,\beta$, so interval realization is an LP on any
graph. With two finite resistance choices per uncertain edge it is
NP-complete on one simple cycle, and so is existential unit-capacity
design on that cycle (\cref{thm:a-bound-realization}). These are
existential questions, whereas robust validation tests every scenario.
